\begin{tikzpicture}[>=Latex, line cap=round, line join=round, font=\small]
\tikzset{
  original/.style={draw=gray!65, dashed, thick},
  general/.style={draw=blue!70!black, fill=blue!9, thick},
  strain/.style={draw=green!45!black, fill=green!10, thick},
  rotation/.style={draw=orange!85!black, fill=orange!10, thick}
}

\begin{scope}
\node[font=\small\bfseries] at (1.7,3.65) {(a) Gradiente local};
\draw[original] (0,0) rectangle (2.5,2.5);
\draw[general] (0.25,-0.1) -- (2.85,0.45) -- (2.35,3) -- (-0.25,2.45) -- cycle;
\draw[->, blue!70!black] (0,0) -- (0.25,-0.1);
\draw[->, blue!70!black] (2.5,0) -- (2.85,0.45);
\draw[->, blue!70!black] (2.5,2.5) -- (2.35,3);
\draw[->, blue!70!black] (0,2.5) -- (-0.25,2.45);
\node[align=center, font=\footnotesize] at (1.3,-0.75)
  {$\{\nabla_j u_i\}$ contiene\\deformación y rotación.};
\end{scope}

\node[font=\Large] at (3.75,1.35) {$=$};

\begin{scope}[xshift=4.65cm]
\node[font=\small\bfseries] at (1.35,3.65) {(b) Parte simétrica};
\draw[original] (0,0) rectangle (2.5,2.5);
\draw[strain] (0,0) -- (2.75,0.35) -- (3.1,3.1) -- (0.35,2.75) -- cycle;
\draw[<->, green!45!black, thick] (0.2,1.15) -- (2.9,1.95);
\draw[<->, green!45!black, thick] (1.15,0.2) -- (1.95,2.9);
\node[green!35!black, font=\large] at (1.55,1.5) {$\mathbf{U}$};
\node[align=center, font=\footnotesize] at (1.5,-0.75)
  {Cambia longitudes\\y ángulos.};
\end{scope}

\node[font=\Large] at (8.45,1.35) {$+$};

\begin{scope}[xshift=9.35cm]
\node[font=\small\bfseries] at (1.25,3.65) {(c) Parte antisimétrica};
\draw[original] (0,0) rectangle (2.5,2.5);
\begin{scope}[rotate around={12:(1.25,1.25)}]
\draw[rotation] (0,0) rectangle (2.5,2.5);
\end{scope}
\draw[->, orange!85!black, very thick] (1.9,1.25) arc[start angle=0,end angle=70,radius=0.65];
\node[orange!85!black, font=\large] at (1.25,1.25) {$\boldsymbol{\Omega}$};
\node[align=center, font=\footnotesize] at (1.25,-0.75)
  {Rotación rígida local: no cambia\\longitudes ni ángulos a primer orden.};
\end{scope}

\node[align=center] at (6.4,-2)
  {$\{\nabla_j u_i\}=\mathbf{U}+\boldsymbol{\Omega}$,\qquad
   $U_{ij}=\dfrac{\nabla_j u_i+\nabla_i u_j}{2}$,\qquad
   $\Omega_{ij}=\dfrac{\nabla_j u_i-\nabla_i u_j}{2}$.};
\node[align=center, font=\footnotesize] at (6.4,-2.75)
  {La parte simétrica $\mathbf{U}$ es el tensor de deformación infinitesimal de Cauchy.};
\end{tikzpicture}